% TOP_PROBLEM: {"id": 3333, "title": "Which lattices occur as intervals in subgroup lattices of finite groups?", "category_id": 17, "category": {"id": 17, "name": "group_theory", "display_name": "Group Theory", "description": "Problems about groups, group actions, representations, and related algebraic structures.", "slug": "group-theory", "order_index": 17, "created_at": "2026-07-31T15:26:25.671Z"}, "status": "open", "problem_number": "OPG-37302", "set_id": 12}
% FIRST_POSTED: 2026-10-06
% ATTEMPT_STATUS: solved
% OPENAI_CLAIM: {"schema_version":1,"source_repo":"https://github.com/openai/math","source_commit":"adc7f1241b42e322a6451854ab7e4b4c146bf78a","release_date":"2026-10-06","match":"full","resolution":"proved","independently_reviewed":false,"second_pass":"agreed","families":[{"family":"206","title":"Finite lattice representation and undecidability","subject":"Algebra","manuscripts":[{"title":"A negative solution to the finite lattice representation problem","pdf_path":"preprints/A-Negative-Solution-to-the-Finite-Lattice-Representation-Problem-September-24-2026/paper.pdf","readme_path":"preprints/A-Negative-Solution-to-the-Finite-Lattice-Representation-Problem-September-24-2026/README.md","date":"2026-09-24","theorem_ref":"main:congruences (sections/introduction.tex:22); main:group (sections/introduction.tex:41)","statement_sources":[{"path":"preprints/A-Negative-Solution-to-the-Finite-Lattice-Representation-Problem-September-24-2026/build/sections/introduction.tex","line":22,"label":"main:congruences"},{"path":"preprints/A-Negative-Solution-to-the-Finite-Lattice-Representation-Problem-September-24-2026/build/sections/introduction.tex","line":41,"label":"main:group"}]},{"title":"Finite congruence lattices: characterization and undecidability","pdf_path":"preprints/Finite-Congruence-Lattices-Characterization-and-Undecidability-September-24-2026/paper.pdf","readme_path":"preprints/Finite-Congruence-Lattices-Characterization-and-Undecidability-September-24-2026/README.md","date":"2026-09-24","theorem_ref":"m:main (sections/introduction.tex:60); n:undecidable (sections/diophantine.tex:154)","statement_sources":[{"path":"preprints/Finite-Congruence-Lattices-Characterization-and-Undecidability-September-24-2026/build/sections/introduction.tex","line":60,"label":"m:main"},{"path":"preprints/Finite-Congruence-Lattices-Characterization-and-Undecidability-September-24-2026/build/sections/diophantine.tex","line":154,"label":"n:undecidable"}]}],"lean":{"doc_path":"lean/docs/206.md","comparators":["lean/ComparatorChallenges/FiniteCongruenceGraph.json","lean/ComparatorChallenges/FiniteCongruenceGraph.lean"],"declarations":[],"covers_main_theorem":false},"reasoning_trace":null}]}
\documentclass[11pt]{article}
\usepackage[margin=25mm]{geometry}
\usepackage{fontspec}
\setmainfont{Latin Modern Roman}
\usepackage{amsmath,amssymb,mathtools}
\usepackage{unicode-math}
\setmathfont{Latin Modern Math}
\usepackage{enumitem}
\usepackage{xurl,hyperref,bookmark}
\hypersetup{colorlinks=true,urlcolor=blue}
\setcounter{secnumdepth}{0}
\setlength{\parindent}{0pt}
\setlength{\parskip}{5pt}
\setlength{\emergencystretch}{3em}
\begin{document}
\section*{Which lattices occur as intervals in subgroup lattices of finite groups?}

\subsection{OpenAI's claim}
The linked manuscripts present the following claims by OpenAI: Some finite lattice is claimed to have no realization as the complete congruence lattice of any finite algebra; A colored-graph criterion is claimed to characterize realizable finite congruence lattices, with an undecidable recognition problem; recognition of finite subgroup intervals is also claimed to be undecidable. Together they give a claimed affirmative answer in the matching scope.

\subsection{Mathematical statement of the claimed result}
\paragraph{Result 1.} The two existence claims concern finite lattices. First, some nonempty finite lattice \(L\) satisfies
\[
 \forall A\text{ finite nonempty algebra of finite signature},\qquad
 L\not\cong\operatorname{Con}(A).
\]
Second, there is a finite lattice \(L'\) for which
\[
 \forall G\text{ finite},\ \forall D\le G,\qquad L'\not\cong[D,G].
\]
Here \(\operatorname{Con}(A)\) is the congruence lattice and \([D,G]\) is the complete interval of subgroups containing \(D\).

\paragraph{Result 2.} For a nonempty finite lattice \(L\), let \(\mathcal R(L)\) mean that \(L\cong\operatorname{Con}(A)\) for some finite algebra \(A\). The claimed finite colored-graph criterion is equivalent to \(\mathcal R(L)\), with a finite decision procedure once the candidate graph size is fixed. Nevertheless,
\[
 \text{no total algorithm decides }\mathcal R(L)\text{ from the order table of }L.
\]
The source also specifies a nonrepresentable lattice \(L_0\) of least cardinality, with \(|L_0|>7\) and every smaller finite lattice representable. Its group-interval assertion is
\[
 \text{no total algorithm decides }\exists G\text{ finite},\ H\le G:\ L\cong[H,G].
\]
The undecidability theorem additionally rules out a total computable bound, depending only on \(|L|\), on the smallest carrier size among representable lattices.

\subsection{Outline of the argument}
The obstruction is first phrased for subgroup intervals using specially tested lattice chains and coordinate transport. Classification-based reductions control possible factor actions, while competing rank and growth estimates contradict the required long configuration; the subgroup-interval equivalence transfers the obstruction to finite algebras.

Finite colored graphs encode the equivalences permitted by all operations of an algebra. Small witness configurations are assembled into lattice representation constraints, and arithmetic divisor obstructions rule out bounded algebraic realizations; a reduction turns the characterization into the claimed undecidability result.

\subsection{Scope relative to this problem}
Theorem 1.2 supplies a finite lattice excluded from every full subgroup interval [H,G] of every finite group, with no normality or group-order restriction, exactly proving the record’s assertion. Theorem 1.1(iv) proves subgroup-interval representability undecidable, which rules out every lattice being representable because that would give the constant yes algorithm, hence proves the record’s negative existence assertion. Within these assumptions, this is a claimed affirmative answer to the stated question.

\subsection{Formal verification}
For family 206, the documented theorem identifies representability with a finite colored-graph witness. Undecidability, nonrepresentation, and subgroup-interval conclusions are not contained in that selected graph equivalence. The listed Comparator files are \texttt{FiniteCongruenceGraph.lean}. The inventory does not supply a named declaration for these Comparator entries. This site has not executed the Lean code. The scope descriptions do not certify the remaining manuscript arguments.

\subsection{Status caveat}
OpenAI attributes these manuscripts to work by an internal model. Its release cautions that arguments without formalization can contain mistakes. This catalogue has not independently checked the proofs; the labels report the mathematical claims attributed to OpenAI.

\subsection{Sources}
\begin{itemize}
\item OpenAI, \emph{A negative solution to the finite lattice representation problem}, 2026-09-24; citation key \texttt{OAI:A-Negative-Solution-to-the-Finite-Lattice-Representation-Problem-September-24-2026}. \href{https://github.com/openai/math/blob/adc7f1241b42e322a6451854ab7e4b4c146bf78a/preprints/A-Negative-Solution-to-the-Finite-Lattice-Representation-Problem-September-24-2026/paper.pdf}{PDF (GitHub)}; \href{https://github.com/openai/math/raw/adc7f1241b42e322a6451854ab7e4b4c146bf78a/preprints/A-Negative-Solution-to-the-Finite-Lattice-Representation-Problem-September-24-2026/paper.pdf}{PDF (direct)}; \href{https://github.com/openai/math/blob/adc7f1241b42e322a6451854ab7e4b4c146bf78a/preprints/A-Negative-Solution-to-the-Finite-Lattice-Representation-Problem-September-24-2026/README.md}{Manuscript page and citation}; \href{https://github.com/openai/math/blob/main/preprints/A-Negative-Solution-to-the-Finite-Lattice-Representation-Problem-September-24-2026/paper.pdf}{Latest version}; Statement source: \href{https://github.com/openai/math/blob/adc7f1241b42e322a6451854ab7e4b4c146bf78a/preprints/A-Negative-Solution-to-the-Finite-Lattice-Representation-Problem-September-24-2026/build/sections/introduction.tex\#L22}{main:congruences}; Statement source: \href{https://github.com/openai/math/blob/adc7f1241b42e322a6451854ab7e4b4c146bf78a/preprints/A-Negative-Solution-to-the-Finite-Lattice-Representation-Problem-September-24-2026/build/sections/introduction.tex\#L41}{main:group}.
\item OpenAI, \emph{Finite congruence lattices: characterization and undecidability}, 2026-09-24; citation key \texttt{OAI:Finite-Congruence-Lattices-Characterization-and-Undecidability-September-24-2026}. \href{https://github.com/openai/math/blob/adc7f1241b42e322a6451854ab7e4b4c146bf78a/preprints/Finite-Congruence-Lattices-Characterization-and-Undecidability-September-24-2026/paper.pdf}{PDF (GitHub)}; \href{https://github.com/openai/math/raw/adc7f1241b42e322a6451854ab7e4b4c146bf78a/preprints/Finite-Congruence-Lattices-Characterization-and-Undecidability-September-24-2026/paper.pdf}{PDF (direct)}; \href{https://github.com/openai/math/blob/adc7f1241b42e322a6451854ab7e4b4c146bf78a/preprints/Finite-Congruence-Lattices-Characterization-and-Undecidability-September-24-2026/README.md}{Manuscript page and citation}; \href{https://github.com/openai/math/blob/main/preprints/Finite-Congruence-Lattices-Characterization-and-Undecidability-September-24-2026/paper.pdf}{Latest version}; Statement source: \href{https://github.com/openai/math/blob/adc7f1241b42e322a6451854ab7e4b4c146bf78a/preprints/Finite-Congruence-Lattices-Characterization-and-Undecidability-September-24-2026/build/sections/introduction.tex\#L60}{m:main}; Statement source: \href{https://github.com/openai/math/blob/adc7f1241b42e322a6451854ab7e4b4c146bf78a/preprints/Finite-Congruence-Lattices-Characterization-and-Undecidability-September-24-2026/build/sections/diophantine.tex\#L154}{n:undecidable}.
\item Family 206: \href{https://github.com/openai/math/blob/adc7f1241b42e322a6451854ab7e4b4c146bf78a/lean/docs/206.md}{Lean scope}.
\item \href{https://github.com/openai/math/blob/adc7f1241b42e322a6451854ab7e4b4c146bf78a/lean/ComparatorChallenges/FiniteCongruenceGraph.json}{Comparator FiniteCongruenceGraph.json}.
\item \href{https://github.com/openai/math/blob/adc7f1241b42e322a6451854ab7e4b4c146bf78a/lean/ComparatorChallenges/FiniteCongruenceGraph.lean}{Comparator FiniteCongruenceGraph.lean}.
\item \href{https://github.com/openai/math/blob/adc7f1241b42e322a6451854ab7e4b4c146bf78a/CONTENTS.md}{OpenAI catalogue entry}.
\end{itemize}
\end{document}
